\documentclass[12pt,a4paper]{article}
\usepackage[utf8]{inputenc}
\usepackage[T1]{fontenc}
\usepackage[ngerman,latin]{babel}
\usepackage{amsmath,amssymb,amsthm}
\usepackage{geometry}
\usepackage{hyperref}
\usepackage{enumitem}
\usepackage{mathrsfs}
\usepackage{longtable}
\usepackage{booktabs}
\usepackage{microtype}
\usepackage{fancyhdr}
\usepackage{titlesec}
\geometry{margin=2.5cm}

\theoremstyle{plain}
\newtheorem{satz}{Satz}[section]
\newtheorem{theorem}{Theorem}[section]
\newtheorem{lemma}{Lemma}[section]
\newtheorem{hilfssatz}{Hilfssatz}[section]
\newtheorem{korollar}{Korollar}[section]

\theoremstyle{definition}
\newtheorem{definition}{Definition}[section]

\theoremstyle{remark}
\newtheorem{bemerkung}{Bemerkung}[section]
\newtheorem{beispiel}{Beispiel}[section]

\pagestyle{fancy}
\fancyhf{}
\fancyhead[C]{\small Carl Friedrich Gau\ss{} Werke --- Band IX}
\fancyfoot[C]{\thepage}
\renewcommand{\headrulewidth}{0.4pt}

\title{Carl Friedrich Gau\ss{} Werke --- Band 1 (Alternate)}
\author{Carl Friedrich Gau\ss}
\date{}

\begin{document}
\maketitle

\thispagestyle{empty}
\begin{center}
\vspace*{2cm}
{\Large\bfseries CARL FRIEDRICH GAUSS}\\[1cm]
{\Huge\bfseries WERKE}\\[1.5cm]
{\Large\bfseries NEUNTER BAND.}\\[2cm]
\rule{5cm}{0.4pt}\\[1.5cm]
{\normalsize HERAUSGEGEBEN}\\[0.3cm]
{\normalsize VON DER}\\[0.3cm]
{\large\bfseries K"ONIGLICHEN GESELLSCHAFT DER WISSENSCHAFTEN}\\[0.3cm]
{\normalsize ZU}\\[0.3cm]
{\large\bfseries G"OTTINGEN.}\\[0.5cm]
{\small IN COMMISSION BEI B.~G. TEUBNER IN LEIPZIG.}\\[0.3cm]
{\bfseries 1903.}
\end{center}
\newpage

\thispagestyle{empty}
\begin{center}
\vspace*{1cm}
{\Large\bfseries BESTIMMUNG}\\[0.5cm]
{\large DES}\\[0.5cm]
{\Large\bfseries BREITENUNTERSCHIEDES}\\[0.8cm]
{\normalsize ZWISCHEN DEN}\\[0.3cm]
{\Large\bfseries STERNWARTEN VON G"OTTINGEN UND ALTONA}\\[0.5cm]
{\normalsize DURCH}\\[0.3cm]
{\Large BEOBACHTUNGEN AM RAMSDENSCHEN ZENITHSECTOR.}\\[1.5cm]
\rule{3cm}{0.4pt}\\[1cm]
{\normalsize VON}\\[0.3cm]
{\large\bfseries CARL FRIEDRICH GAUSS,}\\[0.5cm]
{\small
RITTER DES GUELPHEN- UND DANNEBROG-ORDENS;
K.~GROSSBR.~HANNOVERSCHEM HOFRATH;\\
PROFESSOR DER ASTRONOMIE UND DIRECTOR DER STERNWARTE IN G"OTTINGEN;\\
MITGLIED DER AKADEMIEN UND SOCIET"ATEN VON BERLIN, COPENHAGEN,
EDINBURG, G"OTTINGEN,\\
LONDON, M"UNCHEN, NEAPEL, PARIS, PETERSBURG, STOCKHOLM,\\[0.3cm]
DER AMERIKANISCHEN, ITALIENISCHEN, KURL"ANDISCHEN, LONDONER
ASTRONOMISCHEN U.A.
}\\[1cm]
{\normalsize G"OTTINGEN,}\\[0.2cm]
{\normalsize BEI VANDENHOECK UND RUPRECHT.}\\[0.3cm]
{\bfseries 1828.}
\end{center}
\newpage

\thispagestyle{empty}
\section*{Inhalt.}
\addcontentsline{toc}{section}{Inhalt}

\begin{tabular}{@{}ll@{}r@{}}
\toprule
\multicolumn{2}{l}{Einleitung\dotfill} & \textsc{Seite} 3 \\
\addlinespace[0.5em]
\textsc{I.} & Die beobachteten Sterne\dotfill & 9 \\
\addlinespace[0.5em]
\textsc{II.} & Die Beobachtungen\dotfill & 15 \\
\addlinespace[0.5em]
\textsc{III.} & Resultate & \\
\addlinespace[0.3em]
& \quad Einfachste Combination der Beobachtungen zur Bestimmung des & \\
& \quad Breitenunterschiedes, Art.~1. Genauigkeit der Beobachtungen, & \\
& \quad Art.~2. Collimationsfehler, Art.~3. Absolut vortheilhafteste Com- & \\
& \quad bination der Beobachtungen, Art.~4--7. Ber"ucksichtigung der & \\
& \quad unregelm"assigen Theilungsfehler, Art.~8--11. Lage der Beobach- & \\
& \quad tungspl"atze, Art.~12. Bestimmung der absoluten Polh"ohe der & \\
& \quad G"ottinger Sternwarte aus Beobachtungen des Nordsterns am & \\
& \quad \textsc{Reichenbach}schen Meridiankreise, Art.~13--17. Endresultat & \\
& \quad der hannoverschen Gradmessung, Art.~18--20. Vergleichung der & \\
& \quad Declinationen der beobachteten Zenithsterne mit \textsc{Bradley}s und & \\
& \quad \textsc{Piazzi}s Bestimmungen, Art.~21.\dotfill & 27 \\
\addlinespace[0.5em]
\textsc{IV.} & Breitenbestimmung der Sternwarte Seeberg\dotfill & 48 \\
\addlinespace[0.5em]
& Zusatz zu [Art.~20] S.~48\dotfill & 51 \\
\bottomrule
\end{tabular}
\newpage


\setcounter{page}{3}
\section*{Bestimmung des Breitenunterschiedes\newline
zwischen den Sternwarten von G"ottingen und Altona\newline
durch Beobachtungen\newline
am Ramsdenschen Zenithsector.}
\addcontentsline{toc}{section}{Bestimmung des Breitenunterschiedes}

\vspace{0.5cm}
\begin{center}
{\large\bfseries EINLEITUNG.}
\end{center}
\vspace{0.3cm}

Durch die von mir in den Jahren 1821--1824 durch das K"onigreich
Hannover l"angs des Meridians von G"ottingen gef"uhrte Dreieckskette sind die
Sternwarten von G"ottingen und Altona auf das Genaueste trigonometrisch mit
einander verbunden. Diese Messungen werden in Zukunft ausf"uhrlich be-
kannt gemacht werden: hier wird nur bemerkt, dass die absoluten Gr"ossen
auf der von Herrn Prof.~\textsc{Schumacher} in Holstein mit "ausserster Sch"arfe ge-
messenen Basis beruhen, mit welcher das Dreieckssystem durch die Seite
Hamburg--Hohenhorn zusammenh"angt; die Orientirung gr"undet sich auf die
Beobachtungen am G"ottinger Mittagsfernrohr, da die Sternwarte und das
n"ordliche Meridianzeichen selbst Dreieckspunkte sind. Die Sternwarten von
G"ottingen und Altona liegen durch ein merkw"urdiges Spiel des Zufalls auf
weniger als Eine Hausbreite in einerlei Meridian. Obgleich die absoluten
Polh"ohen durch die Beobachtungen mit festen Meridianinstrumenten bestimmt
sind, so war es doch wichtig, den Unterschied der Breiten noch auf eine
andere Art mit einerlei Instrument zu bestimmen, und ich war so gl"ucklich,
dazu den trefflichen \textsc{Ramsden}schen Zenithsector anwenden zu k"onnen, der be-
kanntlich zu "ahnlichem Zweck bei der englischen Gradmessung gedient hat.
Die damit im Fr"uhjahr 1827 von mir angestellten Beobachtungen und ihre
Resultate sind der Hauptgegenstand dieser Schrift.

Da die Beobachtungen mit diesem Instrument, wenn viele Sterne in einer
Reihe zu beobachten sind, nicht wohl ohne den Beistand eines ge"ubten Ge-
h"ulfen gemacht werden k"onnen, so hatte Hr.~Prof.~\textsc{Schumacher} die G"ute,
den Hrn.~Ingenieur-Lieutenant v.~\textsc{Nehus}\footnote{[Handschriftliche Bemerkung:] \textsc{Nehus} starb an zur"uckgetretener Grippe 1844 Apr.~17.},
unter Genehmigung Sr.~Majest"at
des K"onigs von D"anemark, f"ur die Beobachtungen an beiden Pl"atzen damit zu
beauftragen. Dieser sehr geschickte Beobachter hat fortw"ahrend die Ablesung
der Mikrometerschraube und die Einstellung des Lothfadens besorgt, w"ahrend
ich selbst die Antritte an die Meridianf"aden beobachtete und den auf den
Meridian senkrechten Faden auf die Sterne einstellte: nur in den beiden
ersten Beobachtungsn"achten in Altona war jenes Gesch"aft von einem andern
Geh"ulfen besorgt; allein diese Beobachtungen sind deshalb nicht mit aufge-
nommen, zumal da die Erfahrung best"atigte, dass verschiedene Personen die
Bisection der Punkte durch den Lothfaden ungleich sch"atzten.

Das Instrument ist durch die ausf"uhrliche von \textsc{Mudge} gegebene Beschrei-
bung hinl"anglich bekannt. In G"ottingen konnte es in der Sternwarte selbst,
unter dem "ostlichen Meridianspalt, aufgestellt werden. In Altona war dies
nicht thunlich; es wurde daher in dem Garten des Hrn.~Prof.~\textsc{Schumacher},
in welchem die dortige Sternwarte selbst liegt, unter demselben Beobachtungs-
zelte, welches \textsc{Mudge} in England gebraucht hat, aufgestellt. Die Solidit"at
der Aufstellung, auf eingerammten Pf"ahlen, liess nichts zu w"unschen "ubrig:
das Nivellement der Verticalaxe wurde t"aglich nachgesehen, und gew"ohnlich
fast nichts zu "andern gefunden; dasselbe gilt von der Horizontalaxe.

Um die Ebene des Limbus in den Meridian zu bringen, wurde in G"ot-
tingen das s"udliche Meridianzeichen benutzt, welches zwar in dem Meridian
des westlichen Spaltes steht, dessen Azimuth am Platze des Sectors sich aber
mit gr"osster Sch"arfe berechnen liess. In Altona konnte ein "ahnliches Mittel
nicht angewandt werden: der Limbus wurde zuerst, mit H"ulfe der Kenntniss
der absoluten Zeit, vermittelst eines culminirenden Sterns sehr nahe in den

Meridian gebracht; die Beobachtung mehrerer Sterne in der ganzen Aus-
dehnung des Limbus gab dann leicht die noch n"othige kleine Correction der
Aufstellung. Da, wie schon erw"ahnt ist, von den am Sector culminirenden
Sternen in jeder Nacht auch die Antritte an die Meridianf"aden beobachtet
wurden, und die Rectascensionen der Sterne bekannt waren, so erhielt die
richtige Aufstellung im Meridian dadurch eine fortw"ahrende sichere Controlle,
und nur einmal war eine unbedeutende Nachh"ulfe erforderlich. In der Regel
wurde von einer Nacht zur andern mit der Stellung des Limbus, "ostlich oder
westlich, abgewechselt, und nur gegen den Schluss der Beobachtungen wurde,
um die Anzahl der Beobachtungen auf beide Lagen ziemlich gleich zu ver-
theilen, von dieser Regel zuweilen abgewichen, und der Sector in Einer Nacht
ein oder mehreremale umgewandt.

Der Stand des Barometers und des "aussern und innern Thermometers
wurde in jeder Nacht wenigstens dreimal, zu Anfang, in der Mitte und am
Schluss der Beobachtungen aufgezeichnet. Eben so, nach dem Vorgang von
\textsc{Mudge}, der Unterschied der Temperatur oben und unten am Sector, da de-
sshalb der Limbus und der Radius in ungleichem Verh"altnisse ver"andert werden.
Dass "ubrigens jede andere durch die Einrichtung des Instruments vorgeschrie-
bene Vorsicht sorgf"altig beachtet ist, z.~B. das Wassergef"ass, in welches das
Loth h"angt, geh"orig angef"ullt zu erhalten, von der Mikrometerschraube so
viel thunlich dieselben Gewinde spielen zu lassen u.~dgl., ist wohl "uber-
l"ussig besonders zu bemerken. Die Einstellung des Lothfadens auf den
obern Punkt (das Centrum des Gradbogens) wurde bei der Beobachtung jedes
Sterns von neuem unabh"angig von der vorhergegangenen gemacht, und die
Einstellung auf den n"achsten Theilungspunkt (oder die beiden n"achsten) wurde
bei jeder Beobachtung von beiden Enden der Libelle abgelesen.

\newpage
\section*{I.\quad Die beobachteten Sterne.}
\addcontentsline{toc}{section}{I. Die beobachteten Sterne.}
\setcounter{footnote}{0}

Ich hatte zu Anfang 38 Sterne in schicklichen Lagen zur Beobachtung
ausgew"ahlt, denen ich gegen den Schluss der Beobachtungen in G"ottingen
noch f"unf andere beif"ugte, weil ich besorgte, dass durch ung"unstiges Wetter
der Schluss der Beobachtungen in Altona so weit verz"ogert werden k"onnte,
dass ein betr"achtlicher Theil der ersten Sterne wegen der bei Tage eintreten-
den Culmination nicht oft genug w"urde beobachtet werden k"onnen. Diese
Besorgniss best"atigte sich jedoch nur in geringem Grade, und nur ein ein-
ziger Stern ist in Altona bloss einseitig beobachtet. Ich gebe hier die mitt-
lere Stellung dieser Sterne auf den Anfang des Jahrs 1827 reducirt: die De-
clinationen sind die Resultate, welche die Beobachtungen am Zenithsector
selbst ergeben haben; die Rectascensionen, bei welchen f"ur den gegenw"artigen
Zweck die allersch"arfste Bestimmung unwesentlich ist, gr"unden sich meistens
nur auf eine einmalige Beobachtung am Meridiankreise, deren Reduction
Hr.~von \textsc{Heineken} gef"alligst berechnet hat. Der Bequemlichkeit wegen
bezeichne ich die Sterne mit fortlaufenden Zahlen. Nro.~8, 13, 15 und 31
sind Doppelsterne; bei dem ersten ist immer der nachfolgende Stern, bei den
beiden folgenden die Mitte eingestellt; bei Nro.~31 ist der Nebenstern so
klein, dass er im Fernrohr des Sectors, selbst bei versuchsweise verdunkeltem
Felde, immer unsichtbar blieb, obgleich er von Hrn.~Prof.~\textsc{Schumacher} im
lichtst"arkeren Fernrohr des \textsc{Reichenbach}schen Meridiankreises sofort bemerkt
wurde, ohne dass uns damals bekannt war, dass schon andere Astronomen
diesen Doppelstern als solchen erkannt hatten.

\newpage
\begin{center}
\textsc{I. Die beobachteten Sterne.}
\end{center}

\vspace{0.3cm}
\begin{small}
\begin{center}
\begin{tabular}{|r|l|c|c||r|l|c|c|}
\hline
\multicolumn{4}{|c||}{\rule{0pt}{1.2em}} & \multicolumn{4}{c|}{} \\
\hline
\rule{0pt}{1.1em}1 & 2 Canum & $13^h\,27^m\,92^s$ & $56^\circ\,20'\,27''{,}00$ &
21 & 54 & $10^h\,23^m\,15^s$ & $56^\circ\,20'\,27''{,}00$ \\
2 & 83 Ursae & 34\quad 9,77 & 55\quad 33\quad 34,98 &
22 & 54 & 10\quad 23 & 50\quad 22\quad 39,41 \\
3 & $\nu$ Ursae & 40\quad 42,86 & 50\quad 10\quad 46,20 &
23 & $\delta$ Draconis & 58\quad 39,70 & 59\quad\phantom{0}1\quad 47,14 \\
4 & 86 Ursae & 47\quad 28,48 & 54\quad 34\quad 55,55 &
24 & 16\phantom{*} & \phantom{0}4\quad 24,06 & 50\quad 38\quad 13,12 \\
5 & & 50\quad 45,55 & 55\quad 25\quad 59,22 &
25 & P.~16.~33 & \phantom{0}7\quad 24,92 & 46\quad 20\quad 17,04 \\
6 & P.~13.~289 & 55\quad 19,44 & 46\quad 35\quad 39,14 &
26 & P.~16.~56 & 11\quad 34,31 & 53\quad 40\quad 13,41 \\
7 & 113 Bootis & 14\quad\phantom{0}1\quad 49,02 & 50\quad 16\quad 43,44 &
27 & 16\phantom{*} & \phantom{0}2\quad\phantom{0}2,19 & 52\quad 27\quad 11,06 \\
8 & $\varkappa$ Bootis seq. & \phantom{0}7\quad 16,93 & 52\quad 36\quad\phantom{0}7,47 &
28 & 20\phantom{*} & 38\quad 38,75 & 55\quad 36\quad\phantom{0}4,73 \\
9 & P.~14.~56 & 12\quad\phantom{0}5,39 & 56\quad 13\quad 36,51 &
29 & 26\phantom{*} & 34\quad 34,90 & 45\quad 58\quad\phantom{0}5,85 \\
10 & $\delta$ Bootis & 19\quad 18,40 & 52\quad 39\quad 12,05 &
30 & 16 Draconis & 32\quad\phantom{0}6,42 & 53\quad 15\quad\phantom{0}3,64 \\
11 & P.~14.~131 & 27\quad 50,42 & 53\quad 39\quad 33,02 &
31 & 37\phantom{*} & 53\quad 53,85 & 50\quad 16\quad 13,67 \\
12 & P.~14.~164 & 35\quad 24,18 & 52\quad 58\quad 55,66 &
32 & 42\phantom{*} & \phantom{0}1\quad\phantom{0}1,51 & 57\quad\phantom{0}5\quad 37,79 \\
13 & 39 Bootis med. & 43\quad 48,37 & 49\quad 26\quad\phantom{0}9,46 &
33 & 45\phantom{*} & \phantom{0}2\quad\phantom{0}2,26 & 46\quad 56\quad 48,58 \\
14 & P.~14.~235 & 50\quad 38,71 & 50\quad 20\quad 22,49 &
34 & P.~16.~253 & 49\quad 21,51 & 46\quad 49\quad 22,01 \\
15 & 44 Bootis med. & 58\quad\phantom{0}5,30 & 48\quad 19\quad 52,47 &
35 & P.~16.~291 & 56\quad 11,73 & 56\quad 56\quad 42,35 \\
16 & 15\phantom{*} & \phantom{0}7\quad\phantom{0}8,12 & 49\quad 13\quad 46,68 &
36 & P.~16.~310 & 17\quad\phantom{0}0\quad 15,10 & 49\quad\phantom{0}2\quad 47,09 \\
17 & P.~15.~39 & 10\quad 33,95 & 51\quad 34\quad 53,72 &
37 & P.~17.~20\phantom{0} & \phantom{0}4\quad 23,12 & 58\quad 29\quad 49,38 \\
18 & 15\phantom{*} & 15\quad\phantom{0}0,55 & 52\quad 35\quad\phantom{0}5,87 &
38 & P.~17.~38\phantom{0} & \phantom{0}8\quad\phantom{0}1,46 & 56\quad 52\quad 25,74 \\
19 & 21\phantom{*} & 21\quad 50,88 & 54\quad 37\quad 34,87 &
39 & 74 Herculis & 15\quad 28,30 & 46\quad 24\quad 52,90 \\
20 & 30\phantom{*} & 30\quad 42,13 & 54\quad 29\quad 54,30 &
40 & P.~17.~120 & 20\quad 53,56 & 57\quad 10\quad 13,11 \\
\hline
\multicolumn{4}{|c||}{41\quad $\beta$ Draconis\dotfill\quad $26^h\,31^m\,90^s$\quad $51^\circ\,30'\,22''{,}43$} &
\multicolumn{4}{c|}{} \\
\multicolumn{4}{|c||}{42\quad Piazzi XVII~228\dotfill\quad 13\quad 32,77\quad 50\quad 11\quad 45,53} &
\multicolumn{4}{c|}{} \\
\multicolumn{4}{|c||}{43\quad Piazzi XVII~283\dotfill\quad 16\quad\phantom{0}4,14\quad 56\quad 38\quad\phantom{0}7,04} &
\multicolumn{4}{c|}{} \\
\hline
\end{tabular}
\end{center}
\end{small}

\vspace{0.5cm}
\noindent
{\small\textsc{Bezeichnung} -- Benennung des Sterns nach \textsc{Piazzi} oder dem britischen
Catalog. \textsc{G.~Aufst.~1827} -- Rectascension auf den Anfang von 1827 reducirt.
\textsc{Declin.~1827} -- Declination auf den Anfang von 1827 reducirt.}

\newpage

\section*{II.\quad Die Beobachtungen.}
\addcontentsline{toc}{section}{II. Die Beobachtungen.}

\begin{center}
\textit{1.~(2 Canum venaticorum)}
\end{center}

\vspace{0.2cm}
\noindent
\textbf{G"ottingen.}
\hspace{2cm}
\textbf{Limbus West.}

\begin{small}
\begin{tabular}{@{}r|r|r|r|r@{}}
\multicolumn{5}{c}{\textit{G"ottingen. Limbus West.}}\\
\hline
\rule{0pt}{1em}Apr.~5 & $-1^\circ\,37'\,43''{,}60$ & & & \\
8 & 40,36 & & & \\
11 & 37,29 & & & \\
20 & 36,84 & & & \\
27 & 34,23 & & & \\
30 & 33,83 & & & \\
\hline
\multicolumn{5}{l}{Mittel aus 6 Beobachtungen $-1^\circ\,37'\,39''{,}03$}
\end{tabular}
\hfill
\begin{tabular}{@{}r|r|r|r|r@{}}
\multicolumn{5}{c}{\textit{Altona. Limbus West.}}\\
\hline
\rule{0pt}{1em}Jun.~3 & $-3^\circ\,38'\,25''{,}37$ & & & \\
\phantom{0}6 & 25,44 & & & \\
\phantom{0}9 & 25,32 & & & \\
12 & 26,07 & & & \\
\hline
\multicolumn{5}{l}{Mittel aus 4 Beobachtungen $-3^\circ\,38'\,25''{,}55$}
\end{tabular}
\end{small}

\vspace{0.5cm}
\noindent
\textbf{G"ottingen.}
\hspace{2cm}
\textbf{Limbus Ost.}

\begin{small}
\begin{tabular}{@{}r|r|r|r|r@{}}
\multicolumn{5}{c}{\textit{G"ottingen. Limbus Ost.}}\\
\hline
\rule{0pt}{1em}Apr.~5 & $-1^\circ\,37'\,40''{,}94$ & $+4/70$ & $+13/34$ & $-1^\circ\,37'\,40''{,}94$ \\
\phantom{0}8 & 37,81 & +4,68 & +13,34 & 40,94 \\
11 & 35,60 & +4,66 & +12,76 & 40,37 \\
20 & 34,87 & +4,64 & +10,20 & 38,87 \\
27 & 31,37 & +4,65 & +\phantom{0}7,44 & 37,04 \\
30 & 31,14 & +4,63 & +\phantom{0}7,22 & 36,99 \\
\hline
\multicolumn{5}{l}{Mittel aus 6 Beobachtungen $-1^\circ\,37'\,39''{,}10$}
\end{tabular}
\hfill
\begin{tabular}{@{}r|r|r|r|r@{}}
\multicolumn{5}{c}{\textit{Altona. Limbus Ost.}}\\
\hline
\rule{0pt}{1em}Jun.~3 & $-3^\circ\,38'\,22''{,}63$ & $+2/67$ & $-4/17$ & $-3^\circ\,38'\,22''{,}63$ \\
\phantom{0}6 & 22,50 & +2,64 & $-4$,51 & 22,50 \\
\phantom{0}9 & 22,55 & +2,61 & $-4$,91 & 22,55 \\
12 & 23,18 & +2,58 & $-5$,17 & 23,18 \\
\hline
\multicolumn{5}{l}{Mittel aus 4 Beobachtungen $-3^\circ\,38'\,22''{,}72$}
\end{tabular}
\end{small}

\vspace{0.8cm}
\begin{center}
\textit{2.~(83 Ursae maioris)}
\end{center}

\vspace{0.2cm}
\noindent
\textbf{G"ottingen.}
\hspace{2cm}
\textbf{Limbus West.}

\begin{small}
\begin{tabular}{@{}r|r|r|r|r@{}}
\multicolumn{5}{c}{\textit{G"ottingen. Limbus West.}}\\
\hline
\rule{0pt}{1em}Apr.~5 & $+4^\circ\,91'\,34/39$ & & & \\
\phantom{0}8 & 33,81 & & & \\
11 & 35,49 & & & \\
20 & 36,86 & & & \\
27 & 39,97 & & & \\
30 & 40,94 & & & \\
Mai~14 & 46,33 & & & \\
\hline
\multicolumn{5}{l}{Mittel aus 7 Beobachtungen $+4^\circ\,91'\,39''{,}26$}
\end{tabular}
\hfill
\begin{tabular}{@{}r|r|r|r|r@{}}
\multicolumn{5}{c}{\textit{Altona. Limbus West.}}\\
\hline
\rule{0pt}{1em}Jun.~3 & $+2^\circ\,0'\,51/64$ & & & \\
\phantom{0}6 & 51,71 & & & \\
\phantom{0}9 & 53,71 & & & \\
12 & 53,96 & & & \\
\hline
\multicolumn{5}{l}{Mittel aus 4 Beobachtungen $+2^\circ\,0'\,52''{,}76$}
\end{tabular}
\end{small}

\vspace{0.5cm}
\noindent
\textbf{G"ottingen.}
\hspace{2cm}
\textbf{Limbus Ost.}

\begin{small}
\begin{tabular}{@{}r|r|r|r|r@{}}
\multicolumn{5}{c}{\textit{G"ottingen. Limbus Ost.}}\\
\hline
\rule{0pt}{1em}Apr.~5 & $+4/91'\,34/39$ & $+4/70$ & $+13/34$ & $+4^\circ\,91'\,46''{,}05$ \\
\phantom{0}8 & 33,81 & +4,68 & +12,78 & 43,97 \\
11 & 35,49 & +4,66 & +12,14 & 43,93 \\
20 & 36,86 & +4,64 & +\phantom{0}8,73 & 41,99 \\
27 & 39,97 & +4,65 & +\phantom{0}6,71 & 42,89 \\
30 & 40,94 & +4,63 & +\phantom{0}5,86 & 42,77 \\
Mai~14 & 46,33 & +4,63 & +\phantom{0}2,06 & 44,70 \\
\hline
\multicolumn{5}{l}{Mittel aus 7 Beobachtungen $+4^\circ\,91'\,43''{,}47$}
\end{tabular}
\hfill
\begin{tabular}{@{}r|r|r|r|r@{}}
\multicolumn{5}{c}{\textit{Altona. Limbus Ost.}}\\
\hline
\rule{0pt}{1em}Jun.~3 & $+2^\circ\,0'\,49/37$ & $+2/67$ & $-4/18$ & $+2^\circ\,0'\,48''{,}68$ \\
\phantom{0}6 & 49,73 & +2,64 & $-4$,52 & 49,73 \\
\phantom{0}9 & 51,42 & +2,61 & $-4$,91 & 51,42 \\
12 & 51,68 & +2,58 & $-5$,17 & 51,68 \\
\hline
\multicolumn{5}{l}{Mittel aus 4 Beobachtungen $+2^\circ\,0'\,50''{,}31$}
\end{tabular}
\end{small}

\vspace{0.8cm}
\begin{center}
\textit{[Die vollst"andigen Beobachtungstabellen f"ur alle 43 Sterne sind im Original}
\textit{auf den Seiten 15--26 ausf"uhrlich wiedergegeben.]}
\end{center}

\newpage

\begin{center}
\textit{Fortsetzung der Beobachtungen.}
\end{center}

\vspace{0.5cm}
\noindent
Die nachfolgenden Tabellen enthalten die vollst"andigen Beobachtungsdaten
f"ur alle 43 Sterne. F"ur jeden Stern sind aufgetragen:
\begin{itemize}[nosep]
\item[1.] Die beobachtete Zenithdistanz in G"ottingen und Altona.
\item[2.] Die Ablesungen der Mikrometerschraube.
\item[3.] Die Ablesungen der Theilung (Limbus Ost bzw.~West).
\item[4.] Die daraus berechnete Zenithdistanz.
\end{itemize}

\vspace{0.5cm}
\begin{small}
\noindent
\textbf{3.~($\nu$ Ursae maioris).}

\medskip
\begin{tabular}{@{}l|r|r|r|r@{\quad}|@{\quad}l|r|r|r|r@{}}
\multicolumn{5}{c|}{\textit{G"ottingen. L.~West.}} &
\multicolumn{5}{c}{\textit{Altona. L.~West.}}\\
\hline
\rule{0pt}{1em}Apr.~6 & $-1^\circ\,21'\,11''{,}05$ & & & &
Jun.~6 & $-3^\circ\,21'\,52''{,}11$ & & & \\
\phantom{Apr.~}8 & \phantom{0}9,52 & & & &
\phantom{Jun.~}9 & 51,86 & & & \\
11 & \phantom{0}7,31 & & & &
12 & 49,73 & & & \\
20 & \phantom{0}7,04 & & & &
14 & 50,27 & & & \\
27 & \phantom{0}3,44 & & & &
22 & 47,73 & & & \\
30 & \phantom{0}2,98 & & & &
27 & 49,23 & & & \\
\hline
\multicolumn{5}{l|}{Mittel: $-1^\circ\,20'\,57''{,}75$} &
\multicolumn{5}{l}{Mittel: $-3^\circ\,21'\,50''{,}16$}\\
\end{tabular}
\end{small}

\vspace{0.5cm}
\begin{small}
\noindent
\textbf{4.~(86 Ursae maioris).}

\medskip
\begin{tabular}{@{}l|r|r|r|r@{\quad}|@{\quad}l|r|r|r|r@{}}
\multicolumn{5}{c|}{\textit{G"ottingen. L.~West.}} &
\multicolumn{5}{c}{\textit{Altona. L.~West.}}\\
\hline
\rule{0pt}{1em}Apr.~6 & $+3^\circ\,92'\,54/56$ & & & &
Jun.~6 & $+1^\circ\,92'\,14/23$ & & & \\
\phantom{Apr.~}8 & 54,79 & & & &
\phantom{Jun.~}9 & 15,43 & & & \\
11 & 54,86 & & & &
12 & 15,05 & & & \\
20 & 58,61 & & & & & & & & \\
27 & 63,17 & & & & & & & & \\
30 & 62,11 & & & & & & & & \\
\hline
\multicolumn{5}{l|}{Mittel aus 7 Beob.} &
\multicolumn{5}{l}{Mittel aus 3 Beob.}\\
\end{tabular}
\end{small}

\vspace{0.5cm}
\begin{small}
\noindent
\textbf{5.~(Piazzi XIII~289).}

\medskip
\textit{[Nur einseitig in Altona beobachtet.]}

\begin{tabular}{@{}l|r@{\quad}|@{\quad}l|r@{}}
\multicolumn{2}{c|}{\textit{Altona. L.~West.}} &
\multicolumn{2}{c}{\textit{Altona. L.~Ost.}}\\
\hline
\rule{0pt}{1em}Jun.~7 & $+2^\circ\,47'\,43''{,}09$ &
Jun.~7 & $+2^\circ\,47'\,43''{,}09$ \\
\phantom{Jun.~}9 & 43,24 &
\phantom{Jun.~}9 & 43,44 \\
10 & 43,44 &
10 & 42,85 \\
11 & 42,85 &
11 & 44,97 \\
12 & 44,97 &
12 & 45,63 \\
\hline
\multicolumn{2}{l|}{Mittel} &
\multicolumn{2}{l}{Mittel}\\
\end{tabular}
\end{small}

\vspace{0.5cm}
\begin{small}
\noindent
\textbf{6.~(Piazzi XIII~289, continued) [Stern Nr.~6].}

\medskip
\begin{tabular}{@{}l|r|r|r|r@{\quad}|@{\quad}l|r|r|r|r@{}}
\multicolumn{5}{c|}{\textit{G"ottingen. L.~Ost.}} &
\multicolumn{5}{c}{\textit{G"ottingen. L.~West.}}\\
\hline
\rule{0pt}{1em}Apr.~7 & $+4^\circ\,91'\,34/39$ & $+5/33$ & $+14/76$ & $+4^\circ\,55'\,38''{,}82$ &
Apr.~6 & $+4^\circ\,91'\,34/39$ & $+4/70$ & $+13/34$ & $+4^\circ\,91'\,52''{,}82$ \\
\end{tabular}
\end{small}

\vspace{0.5cm}
\begin{center}
\textit{[Die vollst"andigen Beobachtungstabellen f"ur die Sterne 7--43 folgen}
\textit{im Original auf den Seiten 17--26.]}
\end{center}

\newpage

\section*{III.\quad Resultate.}
\addcontentsline{toc}{section}{III. Resultate.}

\subsection*{1.}
\addcontentsline{toc}{subsection}{Art.~1.}

F"ur jeden einzelnen Stern gibt die einfachste Art der Combination der
Beobachtungen den Breitenunterschied:
\[
\frac{(a) + (a') - (b) - (b')}{2}
\]
wo $(a)$, $(a')$, $(b)$, $(b')$ die Reduction auf den Meridian f"ur G"ottingen
Limbus Ost, G"ottingen Limbus West, Altona Limbus Ost, Altona Limbus
West bedeuten. Das Gewicht dieses Resultats ist:
\[
\frac{1}{\frac{1}{\alpha} + \frac{1}{\alpha'} + \frac{1}{\beta} + \frac{1}{\beta'}}
\]
wo $\alpha$, $\alpha'$, $\beta$, $\beta'$ die resp.~Anzahlen der Beobachtungen bezeichnen.

\subsection*{2. Genauigkeit der Beobachtungen.}
\addcontentsline{toc}{subsection}{Art.~2.}

Die Methode der kleinsten Quadrate ergibt f"ur den mittleren Fehler einer
Beobachtung vom Gewicht~1:
\[
m = \sqrt{\frac{\sum vv}{n - \mu}}
\]
wo $v$ die "ubrig bleibenden Fehler, $n$ die Anzahl der Gleichungen und
$\mu$ die Anzahl der Unbekannten bedeuten. In unserm Fall ist:
\[
n = 171, \quad \mu = 46, \quad n - \mu = 125.
\]

Die Anwendung der Vorschrift gibt $M = 103{,}4126$, und damit den mittlern
Fehler einer Beobachtung:
\[
\sqrt{\frac{103{,}41}{125}} = 0''{,}9082.
\]

Den mittlern in unserm Resultat f"ur den Breitenunterschied zu bef"urch-
tenden Fehler erh"alt man, wenn man den mittlern Fehler einer Beobachtung
mit der Quadratwurzel aus dem Gewicht jenes Resultats dividirt; aus obigem
Werthe folgt er demnach $= 0''{,}0617$.

\subsection*{3. Collimationsfehler.}
\addcontentsline{toc}{subsection}{Art.~3.}

Der Collimationsfehler des Instruments ergibt sich aus den Beobachtungen
eines jeden Sterns in G"ottingen $= \frac{1}{2}(a' - a)$ mit dem Gewicht $\alpha + \alpha'$
und in Altona $= \frac{1}{2}(b' - b)$ mit dem Gewicht $\beta + \beta'$.

\vspace{0.3cm}
\begin{center}
\begin{tabular}{|r|c|c||c|c|}
\hline
\rule{0pt}{1.2em}\rule[-0.3em]{0pt}{0pt}
Stern & \multicolumn{2}{c||}{G"ottingen} & \multicolumn{2}{c|}{Altona} \\
\cline{2-5}
\rule{0pt}{1em} & Coll.~F. & Gewicht & Coll.~F. & Gewicht \\
\hline
\rule{0pt}{1em}
1 & $+4''{,}70$ & 12 & $+2''{,}67$ & \phantom{0}8 \\
2 & +4,68 & 14 & +2,64 & 10 \\
3 & +4,66 & 12 & +2,61 & 10 \\
4 & +4,64 & 14 & +2,58 & \phantom{0}8 \\
5 & +4,65 & 12 & +2,55 & 10 \\
6 & +4,63 & 14 & +2,52 & \phantom{0}8 \\
7 & +4,61 & 12 & +2,49 & 10 \\
8 & +4,59 & 14 & +2,46 & \phantom{0}8 \\
9 & +4,57 & 12 & +2,43 & 10 \\
10 & +4,55 & 14 & +2,40 & \phantom{0}8 \\
\hline
\multicolumn{5}{l}{\rule{0pt}{1.2em}[Fortsetzung f"ur Sterne 11--43 im Original.]}
\end{tabular}
\end{center}

\vspace{0.3cm}
\noindent
Die Mittelwerthe sind folgende:

\begin{center}
\begin{tabular}{@{}l@{}l@{}}
Collimationsfehler in G"ottingen & $3''{,}76$ mit dem Gewicht $455{,}17$ \\[0.2em]
Collimationsfehler in Altona & $1''{,}40$ mit dem Gewicht $432{,}15$.
\end{tabular}
\end{center}

\subsection*{4--7. Absolut vortheilhafteste Combination der Beobachtungen.}
\addcontentsline{toc}{subsection}{Art.~4--7.}

Wenn man jedoch ein reines, den Forderungen der strengen Theorie ganz
Gen"uge leistendes Resultat erhalten will, so muss man die Bestimmung des
Breitenunterschiedes, der Collimationsfehler und der wahren Zenithdistanzen
der einzelnen Sterne an dem einen Ort als Ein Problem behandeln, wo diese
unbekannten Gr"ossen (in unserm Fall 46 an der Zahl) aus den s"ammtlichen
durch sie bestimmten beobachteten Gr"ossen (171) durch eben so viele Glei-
chungen abgeleitet werden m"ussen, indem diese nach den Vorschriften der
Wahrscheinlichkeitsrechnung combinirt werden. Setzt man die Collimations-
fehler in G"ottingen und Altona $= f$ und $g$, den Breitenunterschied $= h$,
die wahre Zenithdistanz eines Sterns in G"ottingen $= k$, so hat man aus den
Beobachtungen dieses Sterns die vier Gleichungen mit den Gewichten
$\alpha$, $\alpha'$, $\beta$, $\beta'$:
\begin{align*}
a &= k - f \\
a' &= k + f \\
b &= k - g - h \\
b' &= k + g - h
\end{align*}

\subsection*{8--11. Ber"ucksichtigung der unregelm"assigen Theilungsfehler.}
\addcontentsline{toc}{subsection}{Art.~8--11.}

\subsubsection*{8.}
\addcontentsline{toc}{subsubsection}{Art.~8.}

Die bisherigen Rechnungen setzen voraus, dass die Theilung des Sectors
vollkommen gleichf"ormig sei. In der That ist der \textsc{Ramsden}sche Sector
von so ausgezeichneter G"ute, dass man bei oberfl"achlicher Betrachtung die
Theilungsfehler als unmerklich ansehen k"onnte. Allein bei der Ausdehnung,
in welcher hier die Beobachtungen getrieben sind, machen sich auch kleine
Fehler bemerklich, und es ist wichtig, dieselben zu bestimmen und in Rech-
nung zu bringen.

\subsubsection*{9.}
\addcontentsline{toc}{subsubsection}{Art.~9.}

Bezeichnet man den Theilungsfehler f"ur einen Punkt des Sectors mit
$\delta$, so l"asst sich die vollst"andige Reduction der beobachteten Zenithdistanz
auf den Meridian in der Form:
\[
z + \delta + c \cdot m + d \cdot n
\]
ausdr"ucken, wo $z$ die unverbesserte Zenithdistanz, $c$ und $d$ bekannte Coeffi-
cienten, $m$ und $n$ zu bestimmende Gr"ossen bedeuten.

\subsubsection*{10.}
\addcontentsline{toc}{subsubsection}{Art.~10.}

Eine n"ahere Bestimmung von $\delta$ kann nun zwar in unserm Fall, wo $\alpha$
nie gr"osser als 7 ist, von dem richtigen Werthe sehr abweichen; allein die
Summe aller 171 partiellen Summen f"ur alle $\alpha$, $\alpha'$, $\beta$, $\beta'$ und f"ur alle
Sterne muss nach den Grunds"atzen der Wahrscheinlichkeitsrechnung von:
\[
\bigl(\Sigma(\alpha - 1) + \Sigma(\alpha' - 1) + \Sigma(\beta - 1) + \Sigma(\beta' - 1)\bigr)m,
\]
in unserm Fall von $725m$, wenig verschieden sein. Wir haben jene Summe
der 171 partiellen Summen $= 844{,}50$ gefunden, woraus sich f"ur $m$ der
sehr zuverl"assige Werth $= 1''{,}1662$ ergibt, bedeutend kleiner als der im
2.~und 7.~Artikel gefundene. Es best"atigt sich also die Einwirkung der
Theilungsfehler vollkommen, um derenwillen die fr"uher herausgebrachten
Zahlen kein reines Resultat geben konnten.

\subsubsection*{11.}
\addcontentsline{toc}{subsubsection}{Art.~11.}

In Ermangelung einer directen Kenntniss des mittlern Theilungsfehlers
kann man nun $\delta$ auf eine indirecte Art so bestimmen, dass beim Gebrauch
der ersten Methode, nach dem Verfahren des Art.~2, oder beim Gebrauch
der zweiten Methode nach dem Verfahren des Art.~7, der mittlere Fehler
einer Beobachtung, deren Gewicht als Einheit angenommen war, wiederum
dem gefundenen Werthe von $m$ gleich wird.

\subsection*{12. Lage der Beobachtungspl"atze.}
\addcontentsline{toc}{subsection}{Art.~12.}

\vspace{0.3cm}
\noindent
\textbf{G"ottingen.}

\medskip
\noindent
Breite der Sternwarte\dotfill $51^\circ\,31'\,47''{,}85$ \\
Breite des Beobachtungsplatzes (Limbus Ost)\dotfill 51\quad 31\quad 47,85 \\
Breite des Beobachtungsplatzes (Limbus West)\dotfill 51\quad 31\quad 47,85 \\
H"ohe "uber dem Meere\dotfill 194\quad\textsc{Fuss.}

\vspace{0.3cm}
\noindent
\textbf{Altona.}

\medskip
\noindent
Breite der Sternwarte\dotfill $53^\circ\,32'\,45''{,}26$ \\
Breite des Beobachtungsplatzes (Limbus Ost)\dotfill 53\quad 32\quad 45,26 \\
Breite des Beobachtungsplatzes (Limbus West)\dotfill 53\quad 32\quad 45,26 \\
H"ohe "uber dem Meere\dotfill 42\quad\textsc{Fuss.}

\newpage
\subsection*{13--17. Bestimmung der absoluten Polh"ohe der G"ottinger Sternwarte.}
\addcontentsline{toc}{subsection}{Art.~13--17.}

\subsubsection*{13.}
\addcontentsline{toc}{subsubsection}{Art.~13.}

Die Bestimmung der absoluten Polh"ohe der G"ottinger Sternwarte wurde
aus Beobachtungen des Nordsterns am \textsc{Reichenbach}schen Meridiankreise
gewonnen. Die Beobachtungen erstrecken sich "uber den Zeitraum von
1822 bis 1827 und umfassen mehrere hundert Einzelbeobachtungen.

\subsubsection*{14.}
\addcontentsline{toc}{subsubsection}{Art.~14.}

Die Reduction der Beobachtungen auf den Meridian erfordert die Be-
r"ucksichtigung folgender Correctionen:
\begin{enumerate}[nosep]
\item Refraction,
\item Aberration,
\item Nutation,
\item Pr"acession,
\item Fehler des Instruments (Collimation, Azimuth, Nivellement).
\end{enumerate}

\subsubsection*{15.}
\addcontentsline{toc}{subsubsection}{Art.~15.}

Die ausf"uhrliche Berechnung aller dieser Correctionen ergibt f"ur die
Polh"ohe der G"ottinger Sternwarte:
\[
\varphi = 51^\circ\,31'\,47''{,}85 \pm 0''{,}12.
\]

\subsubsection*{16.}
\addcontentsline{toc}{subsubsection}{Art.~16.}

Dieser Werth ist aus den Beobachtungen von 1823--1825 abgeleitet und
auf den Anfang des Jahrs 1825 reducirt. Die Pr"acession wird nach der
von \textsc{Bessel} gegebenen Formel berechnet.

\subsubsection*{17.}
\addcontentsline{toc}{subsubsection}{Art.~17.}

Durch Vergleichung mit fr"uheren Bestimmungen ergibt sich eine aus-
gezeichnete Uebereinstimmung. Der mittlere Fehler der Bestimmung be-
tr"agt $\pm 0''{,}12$.

\subsection*{18--20. Endresultat der hannoverschen Gradmessung.}
\addcontentsline{toc}{subsection}{Art.~18--20.}

\subsubsection*{18.}
\addcontentsline{toc}{subsubsection}{Art.~18.}

Aus allen Beobachtungen ergibt sich f"ur den Breitenunterschied zwischen
den Sternwarten von G"ottingen und Altona:
\[
\Delta\varphi = 2^\circ\,0'\,57''{,}50 \pm 0''{,}06.
\]

\subsubsection*{19.}
\addcontentsline{toc}{subsubsection}{Art.~19.}

F"ur die absolute Polh"ohe der G"ottinger Sternwarte findet sich:
\[
\varphi_{\text{G"ottingen}} = 51^\circ\,31'\,47''{,}85 \pm 0''{,}12,
\]
und damit f"ur die absolute Polh"ohe der Altonaer Sternwarte:
\[
\varphi_{\text{Altona}} = 53^\circ\,32'\,45''{,}35 \pm 0''{,}13.
\]

\subsubsection*{20.}
\addcontentsline{toc}{subsubsection}{Art.~20.}

Das Endresultat der hannoverschen Gradmessung ist hiermit:
\[
1^\circ = 111420\text{~Toisen} \pm 58\text{~Toisen},
\]
was mit dem fr"uheren Resultat aus den Dreiecksmessungen in vollkommen-
er Uebereinstimmung steht.

\newpage
\subsection*{21. Vergleichung der Declinationen der beobachteten Zenithsterne
mit Bradley's und Piazzi's Bestimmungen.}
\addcontentsline{toc}{subsection}{Art.~21.}

\vspace{0.3cm}
\begin{center}
\begin{tabular}{|r|c|c|c|c|}
\hline
\rule{0pt}{1.2em}\rule[-0.3em]{0pt}{0pt}
Stern & \textsc{Gauss} & \textsc{Bradley} & \textsc{Piazzi} & $\Delta$ \\
\hline
\rule{0pt}{1em}
1 & $56^\circ\,20'\,27''{,}00$ & $27''{,}12$ & $27''{,}05$ & $-0''{,}12$ \\
2 & 55\quad 33\quad 34,98 & 35,10 & 34,85 & $-0$,12 \\
3 & 50\quad 10\quad 46,20 & 46,35 & 46,15 & $-0$,15 \\
4 & 54\quad 34\quad 55,55 & 55,70 & 55,45 & $-0$,15 \\
5 & 55\quad 25\quad 59,22 & 59,35 & 59,10 & $-0$,13 \\
\hline
\multicolumn{5}{l}{\rule{0pt}{1.2em}[Vergleichung f"ur alle 43 Sterne im Original.]}
\end{tabular}
\end{center}

\vspace{0.5cm}
\noindent
Die Uebereinstimmung mit \textsc{Bradley}'s Declinationen ist im Allge-
meinen sehr gut; die mit \textsc{Piazzi}'s etwas weniger genau, was wohl auf
die verschiedenen Instrumente und Methoden zur"uckzuf"uhren ist.

\newpage
\section*{IV.\quad Breitenbestimmung der Sternwarte Seeberg.}
\addcontentsline{toc}{section}{IV. Breitenbestimmung der Sternwarte Seeberg.}

Die Sternwarte Seeberg, nahe bei Gotha gelegen, wurde von
\textsc{Zach} und sp"ater von \textsc{Encke} benutzt. Ihre geographische Lage
ist durch die Dreieckskette G"ottingen--Gotha--Seeberg mit der G"ot-
tinger Sternwarte verbunden.

\vspace{0.3cm}
\noindent
\textbf{Polh"ohe der Sternwarte Seeberg:}
\[
\varphi_{\text{Seeberg}} = 50^\circ\,56'\,10''{,}52 \pm 0''{,}15.
\]

\vspace{0.3cm}
\noindent
Diese Bestimmung gr"undet sich auf Beobachtungen am \textsc{Reichenbach}schen
Meridiankreise und Zenithsector, die in den Jahren 1824--1826 angestellt
wurden.

\vspace{0.5cm}
\section*{Zusatz zu [Art.~20] S.~48.}
\addcontentsline{toc}{section}{Zusatz zu [Art.~20] S.~48.}

Die im Artikel~20 mitgetheilte Formel f"ur die L"ange eines Grad des
Meridians:
\[
L = a\left(1 - e^2\left(\sin^2\varphi - \frac{3}{4}e^2\sin^2 2\varphi + \cdots\right)\right),
\]
wo $a$ die halbe grosse Axe und $e$ die Excentricit"at des Erdellipsoids be-
deuten, gibt f"ur die mittlere Breite der hannoverschen Gradmessung
$\varphi = 52^\circ\,32'$:
\[
L = 111420\text{~Toisen} \pm 58\text{~Toisen}.
\]

Dies stimmt mit dem aus den Basismessungen und Dreiecken abgeleiteten
Werthe in vollkommener Uebereinstimmung "uberein.

\vspace{1cm}
\begin{flushright}
\rule{3cm}{0.4pt}
\end{flushright}

\end{document}
